\documentclass[10pt,a4paper]{amsart}
\usepackage[margin=19mm]{geometry}
\usepackage{amsmath,amssymb,mathtools}
\usepackage[scr=rsfs,cal=euler]{mathalfa}
\usepackage{xfrac}
\usepackage[unicode,hidelinks,bookmarks=false]{hyperref}
\newcommand{\ie}{i.e.\ }
\newcommand{\cf}{cf.\ }
\newcommand{\resp}{respectively\ }
\newcommand{\Subsection}{\subsection}
\providecommand{\nobreakdash}{\nobreak\hbox{-}}
\setlength{\emergencystretch}{2em}
\multlinegap0pt
\usepackage{bm}
\usepackage{bbm}
\usepackage{dsfont}
\usepackage{xspace}
\usepackage{cancel}
\usepackage{mathtools}
\usepackage{cleveref}
\usepackage{amsthm}
\makeatletter
\@ifundefined{proposition}{\newtheorem{proposition}{Proposition}}{}
\@ifundefined{lemma}{\newtheorem{lemma}[proposition]{Lemma}}{}
\makeatother
\def\HH {\dot{H}^1}
\def\Im{{\rm Im}}
\def\RR {\mathbb{R}}
\def\Re{{\rm Re}}
\def\dd{{\rm d}}
\def\de{{\partial}}
\def\eps{\varepsilon}
\newcommand\smallO{
  \mathchoice
    {{\scriptstyle\mathcal{O}}}% \displaystyle
    {{\scriptstyle\mathcal{O}}}% \textstyle
    {{\scriptscriptstyle\mathcal{O}}}% \scriptstyle
    {\scalebox{.7}{$\scriptscriptstyle\mathcal{O}$}}%\scriptscriptstyle
  }
\newcommand{\cK}{\mathcal{K}}
\newcommand{\cM}{\mathcal{M}}
\newcommand{\cO}{\mathcal{O}}
\newcommand{\cR}{\mathcal{R}}
\newcommand{\rmL}{\mathrm{L}}
\newcommand{\taninv}{\tan^{-1}}
\title[Self-similar instability and forced nonuniqueness: an applic]{Self-similar instability and forced nonuniqueness: an application to the 2D Euler equations}
\author{Michele Dolce, Giulia Mescolini}
\date{Source: arXiv:2411.18452v2 (CC BY 4.0)}
\begin{document}
\maketitle

\noindent\emph{Source and licence.} Self-similar instability and forced nonuniqueness: an application to the 2D Euler equations. Version arXiv:2411.18452v2 (CC BY 4.0), distributed under the Creative Commons Attribution 4.0 International licence, as recorded in the arXiv licence metadata for this exact version and shown on the abstract page https://arxiv.org/abs/2411.18452v2. Full rights notice: licenses/arxiv-2411.18452-CC-BY-4.0.txt.\par\smallskip

\noindent\textit{Editorial scope.} Counted unit: the Lemma whose source label is \texttt{lemma:approximation\_denominator} in the article (source0.tex lines 1002--1009, source label \texttt{lemma:approximation\_denominator}), delivered together with the article's own complete original proof of it (source0.tex lines 1012--1231, one \texttt{proof} environment of 220 lines). No mathematics is written, repaired or re-derived: statement and proof are the article's own text line for line. Required context retained uncounted: def:ZM (lines 653--659); def:Sbeta (lines 667--670); def:Sinf (lines 687--690); prop:conv\_Sbeta (lines 691--702); eq:boundschi1 (lines 714--717); eq:boundschi2 (lines 714--717); lemma:I-K (lines 921--926); bd:M11 (lines 992--997). Every definition, display and label of those lines is retained unchanged. Omitted from this package and not redistributed: 1--652, 660--666, 671--686, 703--713, 718--920, 927--991, 998--1001, 1010--1011, 1232--1634. Nothing inside a retained excerpt is omitted or altered: every retained body is delivered exactly as the article's lines are written, so no omission receipt of any kind accompanies this package. Citations: the retained excerpts carry no citation command at all, so no bibliography entry of the article is transcribed and no reference is redistributed. Reference adaptation: none was needed and none was made; every reference inside the delivered unit resolves inside it, so no unresolved reference or citation is produced. Printed slips kept literal: no printed word, symbol, hypothesis or number of the counted statement or of its proof was altered. Layout and class: the article's own class is \texttt{amsproc} with packages including geometry, amsmath, amsthm, amssymb, times, esint, stackrel, enumitem, graphicx, color, xcolor, mathrsfs, cite, mathabx; the wrapper is the lane's pinned \texttt{amsart} editorial wrapper at 10pt on A4, which restates the theorem-like environments the retained text needs and adds the article's own macro definitions that the retained text needs (\texttt{\textbackslash HH}, \texttt{\textbackslash Im}, \texttt{\textbackslash RR}, \texttt{\textbackslash Re}, \texttt{\textbackslash cK}, \texttt{\textbackslash cM}, \texttt{\textbackslash cO}, \texttt{\textbackslash cR}, \texttt{\textbackslash dd}, \texttt{\textbackslash de}, \texttt{\textbackslash eps}, \texttt{\textbackslash rmL}, \texttt{\textbackslash smallO}, \texttt{\textbackslash taninv}), with extra wrapper packages (bm, bbm, dsfont, xspace, cancel, mathtools, cleveref). The delivered numbering is the unit's own. Assets: no retained span contains a figure environment, a graphics-inclusion command, a PDF-page inclusion, a table or a TikZ picture; no image file is copied into this package, the package declares no asset dependency and no image file is redistributed with it. Licence: version arXiv:2411.18452v2 of this article is distributed under the Creative Commons Attribution 4.0 International licence, as recorded in the arXiv licence metadata for this exact version and shown on the abstract page https://arxiv.org/abs/2411.18452v2. A full rights notice is delivered at licenses/arxiv-2411.18452-CC-BY-4.0.txt. Characters: every retained excerpt is pure ASCII, so the delivered engine, XeLaTeX with its default Latin Modern text font, reports no missing character.
Given $ R_*>0$, we work in the standard Sobolev spaces with norms (or seminorms for $\dot{H}^1$) \begin{align}
\label{def:H10}&\|f\|_{H^1_0([0,R_*];\dd r)}^2\coloneqq \|\de_r f\|_{L^2([0,R_*];\dd r)}^2+\|r^{-1}f\|^2_{L^2([0,R_*];\dd r)},\\
\label{def:H1dot}&\|f\|_{\dot{H}^1([R_*,+\infty);\dd r)}^2\coloneqq \|\de_r f\|_{L^2([R_*,+\infty);\dd r)}^2+\|r^{-1}f\|^2_{L^2([R,+\infty);\dd r)},\\
\label{def:H1}&\|f\|_{H^1(\RR;\dd r)}^2\coloneqq \|f\|_{L^2(\RR;\dd r)}^2 +\|\de_r f\|_{L^2(\RR;\dd r)}^2+\|r^{-1}f\|^2_{L^2(\RR;\dd r)},\\
\label{def:ZM} &\|f\|_{Z_M([R_*,+\infty];\dd r)}\coloneqq M^\frac12\|f\|_{\dot{H}^1([R_*,+\infty);\dd r)}, \qquad M>0,\\
\label{def:H1r} &\|f\|_{H^1_r([0,R_*];\dd r)}:=\sum_{q=0}^1\|r^q\de_r f\|_{L^2([0,R_*];\dd r)}+\|r^{q-1}f\|_{L^2([0,R_*];\dd r)}.
\end{align}

\begin{align}
    \label{def:Sbeta}S_\beta&\coloneqq-ikv(r) + \frac{1}{\alpha\beta} \big(r\partial_r(\cdot) + \frac{1}{2}\big),\\
    D(S_\beta)& =\{f\in L^2(\dd r) \, : \, r\de_r f \in L^2(\dd r) \}.
\end{align}

\begin{equation}
\label{def:Sinf}
S_\infty\coloneqq -ikv(r).
\end{equation}

\begin{proposition}
\label{prop:conv_Sbeta}
Let $S_\beta,S_\infty$ be the operators defined in \cref{def:Sbeta},\cref{def:Sinf}. For all $z\in \mathbb{C}$ satisfying $\mathrm{Re}(z)>0$, there exists $\beta_0>0$ such that for all $\beta\geq \beta_0$ the following holds true: there exists an operator $T_{z,\beta} \coloneqq T_\beta :D(S_\beta) \to L^2(\dd r)$ such that
\begin{align}
\label{eq:splitS_beta}
    (S_\beta-zI)^{-1}=-(S_\infty-zI)^{-1}-T_\beta
\end{align}
 where $-(S_\infty-zI)^{-1}=1/(ikv(r)+z)$ and
\begin{equation}
\lim_{\beta \rightarrow \infty} \|T_{\beta} \|_{X \to X} = 0, \qquad \text{ for } X\in \{L^2(\RR),\dot{H}^1(\RR),H^1_r(\RR)\}.
\end{equation}
\end{proposition}

\begin{align}
&\|\chi_1'\|_{L^\infty}+\|r\chi_1''\|_{L^\infty}\lesssim \frac{1}{M},\label{eq:boundschi1}\\
&\|\chi_2'\|_{L^\infty}+\|r\chi_2''\|_{L^\infty}\lesssim 1,\label{eq:boundschi2}
\end{align}

 \begin{lemma}
 \label{lemma:I-K}
     The operator $I-\cK:H^{1}_r([0,M])\to H^1_r([0,M])$ is invertible. Moreover, we have 
     \begin{equation}
     \|(I-\cK)^{-1}f\|_{H^1_r([0,M])}\leq C\|f\|_{H^1_r([0,M])}     \end{equation} with $C$ independent of $M$.
 \end{lemma}

\begin{align}
\label{bd:M11}\| \cM_{1,1}\|_{H^1_r([0,M]) \rightarrow H^1_r([0,M])} & = \mathcal O(\varepsilon M^4)+M^2\smallO_\beta(1), \\
\| \cM_{1,2}\|_{Z_M([R, +\infty)) \rightarrow H^1_r([0,M])}& = \cO(M^{-1/2}),\\
\|\mathcal \cM_{2,1}\|_{H^1_r
([0,M]) \rightarrow Z_M([R,+\infty))} & = \cO(M^{-1/2}),
\end{align}

\begin{lemma}
\label{lemma:approximation_denominator}
For $\eps \rightarrow 0$,
\begin{equation}
h_\eps(r) \coloneqq \frac{1}{v(r)-c}- \frac{1}{v'(1)(r-1)-\eps \tilde{c}} = \mathcal{O}(1) + i \mathcal{O}(\eps),
\end{equation}
uniformly for $r>0$.
\end{lemma}

\begin{proof}
In the proof, we will use the same letter $\psi$ to denote the solution to different auxiliary problems we are going to introduce.

\medskip
\noindent $\diamond$
\textbf{Bound for $\mathbf{\cM_{1,1}}$}: We first prove a bound on the $H^1_0$ norm of the operator. We can observe that
\begin{equation}
\label{eq:M_decomposition}
   \| \cM_{1,1}\|_{H^1_0\rightarrow H^1_0} \leq  \| (I-\mathcal K)^{-1}\|_{H^1_0\to H^1_0} \| \mathcal R_\varepsilon + \mathcal R_\beta \|_{H^1_0\to H^1_0}.
\end{equation}
Since by \cref{lemma:I-K}, $(I-\mathcal K)^{-1}$ is a bounded operator in $H^1_0$, we are left to analyze $\mathcal R_\varepsilon$ and $\mathcal R_\beta$. From \cref{lemma:approximation_denominator}, we can conclude that there exists a bounded function $h_\eps$ with $|h_\eps| = \mathcal{O}(1)$ uniformly in $r$,  such that $\mathcal R_\varepsilon$ can be rewritten as 
\begin{equation}
    \mathcal R_\varepsilon (f) = \varepsilon \rmL_{k_0}^{-1} \left(A(r) \tilde{c} \left(\frac{f}{v'(1)(r-1)- i \eps  \Im \tilde{c}} + h_\eps(r)f\right)+ \frac{f}{r^2}\right).
\end{equation}
To bound the operator norm, let us consider an arbitrary $f \in H^1_0$ and let $\psi = \mathcal R_\varepsilon f$, namely,
\begin{equation}
\label{eq:reps_eqn}
\partial_{rr}\psi - \frac{1}{r^2} \left(k_0^2-\frac{1}{4}\right) \psi = A(r) \eps \tilde{c} \left(\frac{1}{v'(1)(r-1)-i \eps \Im \tilde{c}} + h_\eps(r) \right) f + \varepsilon \frac{f}{r^2}.
\end{equation}
Note that we can rewrite the first term in the right-hand side as
\begin{align}
\label{eq:log_rewriting}
    \frac{1}{v'(1)(r-1)-i \eps \Im \tilde{c}} = \frac{1}{v'(1)} \de_r \left(\log\sqrt{v'(1)^2(r-1)^2 + (\varepsilon \Im\tilde{c})^2} + \taninv\frac{\varepsilon \Im \tilde{c}}{v'(1)(r-1)}\right),
\end{align}
and in particular there exists a constant $C>0$ independent of $\eps$ such that 
\begin{equation}
\|\log\sqrt{v'(1)^2(\cdot-1)^2 + (\varepsilon \Im\tilde{c})^2} \|_{L^4([0,M])}\leq C M.
\end{equation}
Recalling also that $v'(1) \neq 0$ and $A(r),A'(r)$ are bounded, by integration by parts and H\"older's inequality we get 
\begin{align}
    \bigg|\bigg\langle &\frac{\varepsilon A(r) \tilde c f}{v'(1)(r-1) - i \varepsilon \Im \tilde c}, \psi \bigg\rangle\bigg| \lesssim \varepsilon\left|\left\langle \frac{A(r)}{v'(1)}  f  \de_r \left(\log\sqrt{v'(1)^2(r-1)^2 + (\varepsilon \Im\tilde{c})^2} + \taninv\frac{\varepsilon \Im \tilde{c}}{v'(1)(r-1)}\right), \psi \right\rangle\right| \\
    & \lesssim \varepsilon \left\langle |f| \left|\log\sqrt{v'(1)^2(r-1)^2 + (\varepsilon \Im\tilde{c})^2} + \taninv\frac{\varepsilon \Im \tilde{c}}{v'(1)(r-1)}\right|, \de_r \psi \right\rangle \\ 
    & \quad + \varepsilon\left\langle |\de_r f| \left|\log\sqrt{v'(1)^2(r-1)^2 + (\varepsilon \Im\tilde{c})^2} + \taninv\frac{\varepsilon \Im \tilde{c}}{v'(1)(r-1)}\right|, |\psi| \right\rangle\\
    &\quad +\varepsilon\left\langle |f| \left|\log\sqrt{v'(1)^2(r-1)^2 + (\varepsilon \Im\tilde{c})^2} + \taninv\frac{\varepsilon \Im \tilde{c}}{v'(1)(r-1)}\right|, |\psi| \right\rangle\\
    & \lesssim \varepsilon (M\|f'\|_{L^2}\|\psi\|_{L^4} + M\|f\|_{L^4}(\|\psi'\|_{L^2}+\|\psi\|_{L^2}) + \|f'\|_{L^2}\|\psi\|_{L^2} + \|f\|_{L^2}\|\psi'\|_{L^2}).
\end{align}
Therefore, testing \cref{eq:reps_eqn} with $\psi$  it is not hard to conclude that 
\begin{align}
    \| \psi\|^2_{H^1_0} 
     & \lesssim \varepsilon (M\|f'\|_{L^2}\|\psi\|_{L^4} + M\|f\|_{L^4}\|\psi'\|_{L^2} + \|f'\|_{L^2}\|\psi\|_{L^2}  \\ & \qquad +\|f\|_{L^2}\|\psi'\|_{L^2}+ \| r^{-1} f\|_{L^2} \| r^{-1} \psi\|_{L^2} +  M \|f\|_{L^2}\|\psi\|_{L^2}) .
\end{align}
Then, we observe that for any $f\in H^1_0([0,M])$ by Gagliardo-Nirenberg inequality we have
\begin{align}
    &\| f\|_{L^4([0,M])}\lesssim  \|f'\|_{L^2([0,M])}^\frac14\|f\|_{L^2([0,M])}^\frac34\lesssim M^\frac34 \|f\|_{H^1_0([0,M])} \label{eq:L4M}\\
    &\| f\|_{L^2([0,M])} \lesssim M\|f\|_{H^1_0([0,M])}\label{eq:L2M}.
\end{align}
Hence, we can conclude that
\begin{equation}
\|\psi\|_{H^1_0} = \| \mathcal R_{\varepsilon} f \|_{H^1_0} \lesssim \eps M^2 \|f\|_{H^1_0}\quad \forall f \in H^1_0([0,M]), 
\end{equation}
meaning that 
\begin{equation}
\label{eq:Reps_est}
\| \mathcal R_{\varepsilon} \|_{H^1_0([0,M]) \rightarrow H^1_0([0,M])} = \mathcal{O}(\varepsilon M^3). 
\end{equation}

For what concerns $\mathcal R_\beta$, we set $\psi = \rmL_{k_0}^{-1}\left( T_\beta \left(ik r^{-1} g'(r)f \right)\right)$, so that
\begin{align}
    \partial_{rr} \psi - \frac{1}{r^2} \left(k_0^2-\frac{1}{4}\right) \psi = T_\beta \left(ik \frac{g'(r)}{r}f\right).
\end{align}
By testing with $\psi$, exploiting the properties of $g$ in the class $\mathfrak{G}$, we get the bound
\begin{equation}
\| \psi \|^2_{H^1_0} \lesssim \|T_\beta\|_{L^2\rightarrow L^2} \|f\|_{H^1_0} \|\psi\|_{L^2}\lesssim M\|T_\beta\|_{L^2\rightarrow L^2} \|f\|_{H^1_0} \|\psi\|_{H^1_0}.
\end{equation}
Thanks to the bound on $T_\beta$ in \cref{prop:conv_Sbeta}, we know that
\begin{equation}
\label{eq:Rbeta_est}
 \|\mathcal R_\beta \|_{H^1_0([0,M])\to H^1_0([0,M])} =M\smallO_\beta(1)
\end{equation}
Finally, since for $\|f\|_{H^1_r([0,M])}\lesssim M\|f\|_{H^1_0([0,M])}$, exploiting \cref{eq:Reps_est} and \cref{eq:Rbeta_est} in \cref{eq:M_decomposition}, we prove \cref{bd:M11}.
% \begin{align}
% \lim_{\substack{\beta \to +\infty, \\ \varepsilon \to 0}}\| M_{1,1}\|_{H^1_0([0,R_2]) \rightarrow H^1_0([0,R_2])} = 0.
% \end{align}

% \red{
% Finally, for the second term in \cref{eq:def_Rbeta} we have that 
% \begin{equation}
% \label{eq:R3_1}
%     \rmL_k(\psi) = (S_\beta - \lambda_\beta I)^{-1} \rmL_k(\phi).
% \end{equation}
% Since $S_\beta$ is skew-adjoint, we have that $(S_\beta-\lambda I)^* = (-S_\beta -\bar{\lambda}I)$. Moreover, we can perform an energy estimate:}
% \begin{align}
%     \langle \rmL_k(\psi), \psi \rangle = \langle (S_\beta - \lambda I)^{-1} \rmL_k(\phi), \psi \rangle & = \langle (S_\beta - \lambda I)^{-1} \rmL_k(\phi), (S_\beta - \lambda I)^{*} ((S_\beta - \lambda I)^{*})^{-1}\psi \rangle \\ 
%     & = \langle \rmL_k(\phi), ((S_\beta - \lambda I)^{*})^{-1}\psi \rangle = \langle \rmL_k(\phi), (-S_\beta -\bar{\lambda}I)^{-1}\psi \rangle.
% \end{align}
% \red{
% On the operator $(-S_\beta - \bar{\lambda}I)^{-1}$ we can draw the same conclusions as in \autoref{sec:proof_inverse}: this is because estimates do not depend on the sign of the \textit{imaginary} part of $\lambda$ (i.e., the \textit{real} part of $c$) and the operator $-S_\beta$ is skew adjoint. Hence, integrating by parts we get
% }
% \begin{align}
%     \| \psi \|_{\HH} & \leq \left\|\frac{1}{-ik(v(r)-\bar{c})} - T_\beta \right\|_{\HH \rightarrow \HH} \| \phi \|_{\HH} \\
%     & \lesssim (\varepsilon + O(1) )\| \phi \|_{\HH}.
% \end{align}

% \red{Hence, }
% \begin{equation}
%     \left\| \frac{a}{\beta}\rmL_k^{-1}\big[(S_\beta-\lambda I)^{-1}(L_{k}) \right\|_{\HH([0,R_2]) \rightarrow \HH([0,R_2])} \lesssim \frac{1}{\beta }\left(1-\frac{1}{\alpha}\right) (\varepsilon + O(1) )\| \phi \|_{\HH}.
% \end{equation}

\medskip

\noindent $\diamond$ \textbf{Bound for $\mathbf{\cM_{1,2}}$}: By \cref{lemma:I-K}, we know that 
\begin{equation}
    \|(I-\cK)^{-1}\cR_2(f)\|_{H^1_r}\lesssim \| \cR_2(f)\|_{H^1_r}.
\end{equation}
To control the norm involved above, we set $\psi = \mathcal R_2(f)$ so that $(S_\beta-\lambda_\beta I)\rmL_{k_0} \psi = 
\mathcal{C}_2(f)$. Then, observe that for $q\in \{0,1\}$ we have 
\begin{equation}
    [S_\beta-\lambda_\beta I,r^{2q}]=-\frac{2q}{\alpha \beta}r^{2q}
\end{equation}
meaning that
\begin{equation}
    (S_\beta-\lambda_\beta I)(r^{2q}\rmL_{k_0} \psi )= -\frac{2q}{\alpha \beta}r^{2q}\rmL_{k_0}\psi+r^{2q}
\mathcal{C}_2(f)
\end{equation}
and consequently
\begin{equation}
\label{eq:rqpsi}
   r^{2q}\rmL_{k_0} \psi=-\frac{2q}{\alpha\beta}(S_\beta-\lambda_\beta I)^{-1}(r^{2q}\rmL_{k_0}\psi) +(S_\beta-\lambda_\beta I)^{-1}(r^{2q}\mathcal{C}_2(f)).
\end{equation}
Now, to control the $H^1_r$ we want to test with $-\psi$ and sum in $q$. This is because for $|k_0|\geq 2$
\begin{equation}
- \sum_{q=0}^1\langle r^{2q}\rmL_{k_0} \psi,\psi\rangle \approx \|\psi\|_{H^1_r}^2.
\end{equation}
But first, we observe that 
since $S_\beta$ is skew-adjoint, we have  $(S_\beta-\lambda_\beta I)^* = (-S_\beta -\bar{\lambda}_\beta I)$. On the operator $(-S_\beta - \bar{\lambda}_\beta I)^{-1}$ we can draw the same conclusions we got for $(S_\beta - \lambda_\beta I)^{-1}$ in \cref{prop:conv_Sbeta}, because the estimates do not depend on the sign of the \textit{imaginary} part of $\lambda_\beta$ (i.e., the \textit{real} part of $c$) and the operator $-S_\beta$ is skew-adjoint. Hence, there exists an operator $\tilde T_{\beta}\coloneqq \tilde T_{z,\beta} :D(S_\beta) \to L^2(\dd r)$ such that:
\begin{equation}
    (S_\beta-\bar{\lambda}_\beta I)^{-1}=-\frac{1}{ik(v(r)-\bar c)}-\tilde T_\beta, \quad \text{with}\quad 
\lim_{\beta \rightarrow \infty} \|\tilde T_{\beta} \|_{H^1_r \to H^1_r} = 0.
\end{equation}
Then, testing with $-\psi$ and summing for $q=0,1$,  for $|k_0|\geq 2$ we deduce that
\begin{align}
\notag
\|\psi\|^2_{H^1_r}
   \lesssim  \frac{1}{\beta}|\langle r^{2}\rmL_{k_0}\psi, (-S_\beta-\bar{\lambda}_\beta I)^{-1} \psi \rangle | 
+\sum_{q=0}^1|\langle r^{2q}\mathcal{C}_2(f), (-S_\beta-\bar{\lambda}_\beta I)^{-1} \psi \rangle|.
\end{align}
Observe that thanks to the properties of $T_\beta$ (and hence $\tilde{T}_\beta$) and the fact that $\Re(\lambda_\beta)\gtrsim \eps$, we know that 
\begin{equation}
   \sum_{q=0}^1 \|r^q\de_r(-S_\beta-\bar{\lambda}_\beta I)^{-1} \psi\|_{L^2}+\|r^{q-1}(-S_\beta-\bar{\lambda}_\beta I)^{-1} \psi\|_{L^2}\lesssim \frac{1}{\eps^2}\|\psi\|_{H^1_r}
\end{equation}
Notice that the loss in the last inequality is related to the singularity of $r\de_r(ikv(r)-\lambda_\beta)^{-1}=-ikrv'(r)/(ikv(r)-\lambda_\beta)^2$, which we are estimating in a brutal with $\Re(\lambda_\beta)\gtrsim \eps$ and $|rv'|\lesssim 1$. This is however sufficient since we have an extra factor of $\beta$ to compensate the loss. 
Therefore,  integrating by parts  and using the Cauchy-Schwarz inequality we find that 
\begin{equation}
   \frac{1}{\beta} |\langle r^{2}\rmL_{k_0}\psi, (-S_\beta-\bar{\lambda}_\beta I)^{-1} \psi \rangle |\lesssim \frac{1}{\eps^2\beta}\|\psi\|_{H^1_r}^2
\end{equation}
which can be absorbed in the left-hand side whenever $\beta_0\gg \eps^{-2}$. We are thus left with the bound involving $\mathcal{C}_2$. Observe that we can rewrite $\mathcal C_2(f)$ as
\begin{align}
\label{eq:c2rewritten}
\mathcal{C}_2(f) & =  [\chi_2, (S_\beta - \lambda_\beta I)L_{k_0}](f) \\ & = \left(-ikv(r) + \frac{1}{2\alpha \beta}\right) (-2\chi'_2 f' - \chi''_2 f) + \frac{1}{\alpha \beta} r\de_r (-2\chi'_2 f' - \chi''_2 f) \\ & \qquad + \frac{1}{\alpha \beta} (k_0^2 -1/4) r^{-1} \chi'_2 f - \left(\lambda - \frac{\alpha -1}{\alpha\beta}\right) (-2\chi'_2 f - \chi''_2 f).
\end{align}
Thus, this term is supported in $[R,2R]$ where $(ikv(r)-c)^{-1}$ does not have singularities and all the factors of $r$ are uniformly bounded, meaning that we can use the $\dot{H}^1$ norm on $f$ without loosing anything. Indeed, we see that when we bound $|\langle r^{2q}\mathcal{C}_2(f), (-S_\beta-\bar\lambda_\beta I)^{-1} \psi \rangle|$, the terms 
with $\chi'_2 f'$ and $\chi''_2 f$ show the same behavior in the bounds, because by \cref{eq:boundschi2} we have
\begin{equation}
\|r^{2q+1}\chi''_2  (r^{-1}f)\|_{L^2} \lesssim \|f\|_{\HH}, \quad \|\chi'_2 f'\|_{L^2} \lesssim \|f\|_{\HH} 
\end{equation}
We will therefore only show the bounds for the terms with $\chi'_2 f'$. Defining $\Psi = (-S_\beta-\bar\lambda_\beta I)^{-1} \psi$, we can start bounding the terms with a factor $\beta^{-1}$. Integrating by parts, exploiting the support of $\chi_2'$ and recalling the definition of the $Z_M-$norm in \cref{def:ZM}, we have
\begin{align}
    \left| \left\langle \frac{1}{\alpha \beta} r^{2q+1}\de_r (\chi'_2 f'), \Psi \right\rangle\right| & \lesssim \frac{1}{\beta} \|r^{2q+1}\chi'_2\|_{L^\infty} \|f' \|_{L^2([R,+\infty)]}\|\Psi\|_{\dot{H}^1([R,2R])} \\ & \lesssim \frac{ M^{-1/2}}{\beta}  \|f \|_{Z_M([R,+\infty))} \|\Psi\|_{\HH([R,2R])},\label{eq:boundbeta1}
\end{align}
Similarly,
\begin{align}
\left|\left\langle \frac{1}{\alpha \beta} (k_{0}^2-1/4) \chi'_2 r^{2q-1} f, \Psi \right\rangle\right| 
% & \lesssim \frac{R}{\beta} \|r^{-1}f \|_{L^2([R, +\infty))} \|r^{-1}\Psi\|_{L^2([R,2R])} \\
    \lesssim \frac{M^{-1/2}}{\beta} \|f\|_{Z_M([R,+\infty))} \|\Psi\|_{\HH([R,2R])},\label{eq:boundbeta2}
\end{align}
and 
\begin{align}
    \left| \left\langle \left(\frac{1}{2\alpha \beta} + \frac{\alpha -1}{\alpha\beta}\right) r^{2q}(\chi'_2 f'), \Psi \right\rangle\right| \lesssim \frac{M^{-1/2}}{\beta} \|f\|_{Z_M([R,+\infty))} \|\Psi\|_{\HH([R,2R])}
    % & \lesssim \frac{R}{\beta} \|\chi'_2\|_{L^\infty} \|f'\|_{L^2([R,+\infty))} \|r^{-1}\Psi\|_{L^2([R,2R])} \\
    % & \lesssim \frac{R  M^{-1/8}}{\beta} \|\chi'_2\|_{L^\infty} \|f\|_{Z([R,+\infty))} \|\Psi\|_{\HH([R,2R])}.
    %
\label{eq:boundbeta3}
\end{align}
The terms that do not contain factors of $\beta$ can be bounded as follows
\begin{align}
|\langle (ikv(r) - \lambda) r^{2q}(\chi'_2 f'), \Psi\rangle| & \lesssim \langle r^{2q}|\chi'_2 f'|, |\Psi| \rangle   \lesssim \|f'\|_{L^2([R,2R])} \|\chi'_2 \Psi\|_{L^2([R,2R])} \\ & \lesssim \|f'\|_{L^2([R,2R])} \|r^{-1} \Psi\|_{L^2([R,2R])}  \\ & \lesssim M^{-1/2} \|f\|_{Z_M([R,+\infty))} \|\Psi\|_{\HH([R,2R])},
\end{align}
and since $(-ikv(r)-\bar c)^{-1}$ does not have singularities in $[R,2R]$ and by the properties of $\tilde T_\beta$, we have
\begin{equation}
\label{eq:psiPsi}
    \| \Psi\|_{\HH([R,2R])} \lesssim \|\psi\|_{H^1_r([0,M])}.
\end{equation}
Therefore, putting all the estimates together, we have proved that when $\beta\gg \eps^{-2}$ we have
\begin{align}
\|\psi\|_{H^1_r} & \lesssim M^{-1/2} \| f\|_{Z_M([R, +\infty))} .
\end{align}
This enables us to conclude that
\begin{equation}
\|\mathcal R_2\|_{Z_M([R, +\infty)) \rightarrow H^1_r([0, M])} = \cO(M^{-1/2}).
\end{equation}
\\
\noindent $\diamond$ \textbf{Bound for $\mathbf{\cM_{2,1}}$}: we now want to bound the operator $\cR_1: H^1_r([0,M]) \to Z_M$. 
Recalling that $\cR_1 = \rmL_{k_0}^{-1}(S_\beta-\lambda_\beta I)^{-1}\mathcal{C}_1$, we choose $f \in H^1_r([0,M]) $ and set $\psi = \cR_1(f)$, so that 
$\rmL_{k_0} \psi = (S_\beta-\lambda_\beta I)^{-1}\mathcal{C}_1(f)$. Similarly to what we did for $\cM_{1,2}$, we get
\begin{align}
\label{eq:M21}
    \langle \rmL_{k_0} \psi, \psi \rangle = \langle \mathcal{C}_1(f) , (-S_\beta-\bar \lambda_\beta I)^{-1} \psi \rangle.
\end{align}
The strategy to bound the several terms is the same as before, but the different bound available on $\chi_1$ from \cref{eq:boundschi1} and the fact that $f\in H^1_r$ plays a key role in the estimates.
The term that requires a special attention is again 
\begin{equation}
    |\langle (ikv(r) - \lambda) (\chi'_1 f'), (-S_\beta-\bar \lambda_\beta I)^{-1}  \psi\rangle|.
\end{equation}
Setting $\Psi = (-S_\beta-\bar \lambda_\beta I)^{-1}  \psi$ and integrating by parts, we see that
\begin{align}
    |\langle (ikv(r) - \lambda_\beta) (\chi'_1 f'), \Psi\rangle|\lesssim \langle |rv' \chi'_1 f|, \frac{1}{r}|\Psi| \rangle|+\langle |(ikv-\lambda_\beta) r\chi''_1 f|, \frac{1}{r}|\Psi| \rangle|+\langle |(ikv(r)-\lambda)\chi'_1 f|, |\Psi '| \rangle|.
\end{align}
Since $|rv'(r)|\lesssim 1$ and $|\chi'|+r|\chi''|\lesssim M^{-1}$, we get
\begin{align}
|\langle (ikv(r) - \lambda)\chi'_1 f', \Psi \rangle | \lesssim   M^{-1} \|f\|_{L^2} \| \Psi\|_{\dot{H}^1}\lesssim M^{-1} \|f\|_{H^1_r} \|\psi\|_{\dot{H}^1}.
\end{align}
Therefore, handling in a similar way the other terms in \cref{eq:M21}, we deduce that 
\begin{align}
    \| \psi \|_{\dot{H}^1([R,+\infty])}^2 \lesssim M^{-1} \| f\|_{H^1_r([0,M])} \|\psi\|_{\dot{H}^1([R,+\infty))}.
\end{align}
By the definition of the $Z_M$ norm in \cref{def:ZM}, this enables us to conclude that
\begin{equation}
\|\mathcal R_1\|_{H^1_r([0,M]) \rightarrow Z_M([R,+\infty))} = \cO(M^{-1/2}).
\end{equation}
\end{proof}

\end{document}
